\section{Introduction}
\label{sec:introduction}

Quality-Diversity (QD) optimization seeks not a single optimal solution, but a diverse collection of high-performing solutions spanning a space of user-defined behavioral characteristics~\cite{pugh2016quality, chatzilygeroudis2021quality}.
This paradigm has proven transformative in robotics, where diverse repertoires enable rapid adaptation to damage~\cite{cully2015robots}; in game design, where varied content improves player experience~\cite{gravina2019procedural}; and in engineering, where exploring the design space reveals unexpected trade-offs~\cite{gaier2018data}.

The dominant QD algorithm, MAP-Elites~\cite{mouret2015illuminating}, discretizes the behavior space into a uniform grid and maintains one elite solution per cell.
Its simplicity and effectiveness have driven widespread adoption, spawning variants that improve performance through advanced emitters~\cite{fontaine2020covariance, colas2020scaling}, learned representations~\cite{paolo2021sparse}, and gradient information~\cite{nilsson2021policy}.

However, MAP-Elites suffers from a fundamental limitation: the \emph{curse of dimensionality}.
A grid with $m$ bins per dimension requires $m^d$ cells for a $d$-dimensional behavior space.
Even with a conservative $m = 3$, a 10-dimensional space requires $3^{10} = 59{,}049$ cells, and 20 dimensions requires over 3.5 billion.
Most cells remain empty, the grid provides poor coverage, and computational resources are wasted on an exponentially sparse structure.
CVT-MAP-Elites~\cite{vassiliades2018using} mitigates this by replacing the grid with Voronoi tessellation from $k$-means centroids, but still requires pre-specifying the number and placement of niches, and its computational cost grows with behavior dimension.

We observe that archive maintenance in QD optimization is fundamentally a \emph{subset selection} problem: from a pool of candidate solutions, select $K$ that are collectively high-quality and diverse.
This is precisely the problem solved by Maximum Marginal Relevance (MMR) in information retrieval~\cite{carbonell1998use}, where documents are selected to be both relevant to a query and different from already-selected results.

We introduce \mmr, which directly applies MMR to QD archive maintenance.
Each candidate is scored by a weighted combination of normalized fitness (relevance) and minimum distance to the current archive (novelty), and a greedy selection process builds the archive one solution at a time.
A lazy evaluation strategy with priority-queue-based staleness tracking achieves $O(K \log K)$ expected time complexity, implemented in Rust for practical efficiency.

Our contributions are:
\begin{enumerate}
    \item \textbf{A novel QD algorithm} that eliminates grid discretization by reformulating archive maintenance as MMR selection, enabling operation in arbitrarily high-dimensional behavior spaces.
    \item \textbf{Theoretical grounding} via the connection to submodular maximization, providing approximation guarantees for the diversity component.
    \item \textbf{An efficient implementation} using lazy greedy evaluation with $O(K \log K)$ expected complexity, achieving ${\sim}50\times$ speedup over naive selection.
    \item \textbf{Empirical validation} showing $12\times$ better archive uniformity than \map at 20D, graceful scaling to 100D, and competitive fitness across benchmarks.
\end{enumerate}
